\section*{Erd\H{o}s Problem 1099: power sums of consecutive divisor ratios}
\subsection*{Statement and scope}
For the ordered divisors $1=d_1<\cdots<d_{\tau(n)}=n$, put
\[
 h_\alpha(n)=\sum_{i=1}^{\tau(n)-1}\left(\frac{d_{i+1}}{d_i}-1\right)^\alpha.
\]
The question is whether $\liminf_{n\to\infty}h_\alpha(n)<\infty$ for every $\alpha>1$. Vose proved this in \emph{Integers with consecutive divisors in small ratio}, J. Number Theory 19 (1984), 233--238, \url{https://doi.org/10.1016/0022-314X(84)90107-0}. We do not reconstruct his sufficient construction. We prove a universal lower bound, clarify the critical exponent, and derive the auxiliary sum-of-ratios consequence. The existence assertion remains unresolved in this attempt.

\subsection*{Refining a multiplicative partition decreases the cost}
If $a<b<c$ are positive reals and $\alpha\ge1$, write $u=b/a-1$, $v=c/b-1$. Then
\[
 c/a-1=u+v+uv\ge u+v,\qquad (u+v)^\alpha\ge u^\alpha+v^\alpha.
\]
Consequently inserting $b$ between $a,c$ cannot increase the sum of powered relative gaps. Repeated insertion proves this for any finite refinement.

Let $m=\lfloor\sqrt n\rfloor$. Refine the divisor list by the positive real numbers
\[
 1,2,\ldots,m,\quad n/m,n/(m-1),\ldots,n.
\]
This list already contains every original divisor: one of $d,n/d$ is at most $\sqrt n$ and is an integer. Delete the repeated central point if $n=m^2$. The ratios in its two halves are $(j+1)/j$, in reverse orders, and its middle ratio is $n/m^2$. Thus
\[
 \boxed{h_\alpha(n)\ge
 2\sum_{j=1}^{m-1}\frac1{j^\alpha}
 +\left(\frac n{m^2}-1\right)^\alpha.}
 \tag{1}
\]
This is a comparison with a refinement by real numbers, not an assertion that all the inserted numbers divide $n$.

In particular, for $\alpha>1$,
\[
 \liminf_{n\to\infty}h_\alpha(n)\ge2\zeta(\alpha).
\]
For $\alpha=1$, (1) grows like $\log n$. For $0<\alpha<1$, the elementary inequality $(e^t-1)^\alpha\ge t$ for all sufficiently small positive $t$, together with $(e^t-1)^\alpha/t$ having a positive infimum on $(0,\infty)$, gives $h_\alpha(n)\gg_\alpha\sum_i\log(d_{i+1}/d_i)=\log n$. Alternatively, the following exact identity immediately handles the endpoint $\alpha=1$:
\[
 \sum_i\log(d_{i+1}/d_i)=\log n.
\]
Thus the restriction $\alpha>1$ is necessary. The positive infimum just used follows from the ratio tending to infinity both as $t\downarrow0$ and as $t\to\infty$, and continuity in between.

\subsection*{A necessary abundance of divisors}
For $\alpha>1$, let $u_i=d_{i+1}/d_i-1$. Since $\log(1+u_i)\le u_i$, H\"older's inequality gives
\[
 \log n\le\sum_i u_i
 \le(\tau(n)-1)^{1-1/\alpha}h_\alpha(n)^{1/\alpha}.
\]
Any sequence with $h_\alpha(n)\le H$ must therefore satisfy
\[
 \tau(n)-1\ge H^{-1/(\alpha-1)}
                 (\log n)^{\alpha/(\alpha-1)}.
 \tag{2}
\]
Prime powers alone cannot work: $h_\alpha(p^a)=a(p-1)^\alpha\to\infty$. Nor does a uniformly bounded largest divisor ratio by itself suffice; it controls each summand but not their number.

\subsection*{The auxiliary sum of ratios}
Carefully retaining the $\tau(n)-1$ gaps, define
\[
 E(n)=\sum_{i=1}^{\tau(n)-1}\frac{d_{i+1}}{d_i}-\tau(n)-\log n.
\]
Then
\[
 E(n)=\sum_i\bigl(u_i-\log(1+u_i)\bigr)-1.
\]
For $u\ge0$,
\[
 0\le u-\log(1+u)=\int_0^u\frac t{1+t}\,dt\le\frac{u^2}{2}.
\]
It follows that
\[
 -1\le E(n)\le\frac12h_2(n)-1.
\]
So a bounded sequence for the main question at $\alpha=2$ indeed gives the finite-liminf assertion about $E(n)$. The formula also prevents an off-by-one error in the unpowered sum.

\subsection*{Assessment}
The lower bound (1), the necessity of $\alpha>1$, the divisor-count obstruction (2), and the auxiliary implication are proved. What is missing is a sequence whose entire logarithmic divisor mesh has uniformly bounded powered cost, including its long central region. Neither numerical examples nor the refinement lower bound supplies such a sequence. Vose's theorem is correctly attributed as the known affirmative answer, not inserted as a proof of our missing construction. Source statement: \url{https://www.erdosproblems.com/1099}.
\subsection*{COMPLETION ESTIMATE}
\textbf{COMPLETION: 40\%}
